\scanpage{10}
\mspage{9}
\begin{center}4 - 26 -\end{center}
\margin{n\textsuperscript{o} 392}
\[
 (8)\qquad \chi(\operatorname{Coker}(f))=\operatorname{div}(\det(\underline{U})).
\]
On peut considérer $L_1$ et $L_2$ comme des réseaux dans $V_1=L_1\otimes_A K$ et $V_2=L_2\otimes_A K$ respectivement, $f$ s'étendant en un $K$-homomorphisme $f_{(K)}:V_1\longrightarrow V_2$. On a alors $(\operatorname{Coker}(f))_{(K)}=\operatorname{Coker}(f_{(K)})$ et dire que $\operatorname{Coker}(f)$ est un $A$-module de torsion signifie que $\operatorname{Coker}(f_{(K)})=0$ ; or, il revient au même de dire que $f_{(K)}$ est surjectif ou que $\det(\underline{U})\ne0$, d'où la première assertion. D'autre part, on peut écrire $f(L_1)=u(L_2)$, où $u$ est un endomorphisme de $L_2$ de déterminant $\det(\underline{U})$ ; comme $\operatorname{Coker}(f)=L_2/u(L_2)$, la formule (8) résulte de (7) et (6).

\uline{Exemple}.- Si $A=Z$, le groupe des diviseurs de $A$ s'identifie au groupe multiplicatif $Q_+^*$ des nombres rationnels $>0$. Pour tout groupe commutatif fini $T$, $\chi(T)$ est l'\uline{ordre} de $T$ ; le corollaire précédent montre que l'ordre du groupe $\operatorname{Coker}(f)$ est égal à la \uline{valeur absolue} de $\det(\underline{U})$.

7. \uline{Classes de diviseurs attachéesaux modules de type fini}.

Les notations et hypothèses sont celles des n\textsuperscript{os} $2$ à $6$. Rappelons que l'on note $C(A)$, (ou simplement $C$), le \uline{groupe des classes de diviseurs} de $A$, quotient de $D(A)$ par le sous-groupe des diviseurs principaux. Pour tout diviseur $d\in D$, nous noterons $c(d)$ sa classe dans $C$.

PROPOSITION 15.- \uline{Soit} $M$ \uline{un} $A$-\uline{module de type fini}. \uline{Il existe un sous-module libre} $L$ \uline{de} $M$ \uline{tel que} $M/L$ \uline{soit un module de torsion}, \uline{et l'élément} $c(\chi(M/L))$ \uline{de} $C$ \uline{ne dépend pas du sous-module libre} $L$ \uline{choisi}.

Posons $S=\unsure{A-\{0\}|A:-\{0\}}$, et soit $V=S^{-1}M=M\otimes_A K$ ; si $n$ est la dimension de $V$ sur $K$, il existe $n$ éléments $e_i$, $(1\le i\le n)$, de $M$ dont les images canoniques dans $V$ forment une base de $V$ ; ces éléments sont évidemment linéairement indépendants dans $M$, donc engendrent un sous-module libre $L$ de $M$ tel que $S^{-1}(M/L)=S^{-1}M/S^{-1}L=0$, de sorte que $M/L$ est un module de torsion.
